% ================================================================
%  test_formatting.tex —— 排版测试文件
%
%  这个文件故意包含多种排版不规范的问题，
%  用于测试 LaTeXcellence 的 formatting（自动排版）模块。
%
%  【注意】编译成 PDF 后看不出问题（LaTeX 会自动处理空白），
%  但 .tex 源码的阅读体验很差。formatting 模块的目标就是
%  自动修正这些问题，让源码更整洁。
%
%  下面用注释标出每个"不规范"之处。
% ================================================================

\documentclass{article}

\usepackage{amsmath}
\usepackage{amssymb}
\usepackage{geometry}

\title{Test Document for Formatting}
\author{Author Name}
\date{\today}

\begin{document}

\maketitle

% ================================================================
%  问题 1：标题前后有多余空格
%  \section{  Title  } 花括号内的首尾空格应去除
% ================================================================
\section{  Introduction  }
%                   ^^       ^^  多余的内部空格

% 问题 2：段落缩进不一致
%   这里缩进了 4 个空格，下一段可能缩进 2 个或 6 个
%   规范做法：顶格写或用统一的缩进
    This is the first paragraph with weird indentation.
    It has   extra spaces   between words.
%          ^^^         ^^^  单词间多余空格（应合并为一个）
     It also has a strange line break here.
%    ^ 此行前多缩进
The next sentence should be on the same paragraph.

% 问题 3：章节之间空行不统一
%   这里只有 1 个空行，其他地方有 2 个甚至 3 个
  \section{Methodology}

% 问题 4：二级标题也有多余空格
\subsection{  Subsection  with   spaces  }

We propose a novel approach using machine learning techniques.
Our method consists of three main steps.

First, we collect data from various sources. Then we preprocess it.
Finally, we apply our algorithm.

% ================================================================
%  问题 5：多个独立的 equation 环境应合并为 align
%  下面三个公式各自独立编号，实际上它们属于同一组
% ================================================================
\begin{equation}
a + b = c
\end{equation}
\begin{equation}
x ^ 2 + y ^ 2 = z ^ 2
\end{equation}
\begin{equation}
E = mc ^ 2
\end{equation}
% 推荐写法：
% \begin{align}
%   a + b &= c \\
%   x^2 + y^2 &= z^2 \\
%   E &= mc^2
% \end{align}
% （用 align 共享一个编号，或者用 align* 不编号）

These equations should ideally be in an align environment.

% 问题 6：equation 内嵌 aligned 写法较冗余
%   如果是单个公式，直接用 equation + split 更简洁
\begin{equation}
  \begin{aligned}
    f(x) &= ax ^ 2 + bx + c \\
    g(x) &= \sin(x) + \cos(x)
  \end{aligned}
\end{equation}

  \subsection{Another subsection}
%^ 不需要的缩进

Some text here.

% ================================================================
%  问题 7：itemize 列表中 \item 缩进混乱
%  每个 \item 前的空格数不同：2 / 7 / 2 / 16 格
%  规范：所有 \item 在同一列对齐
% ================================================================
  \begin{itemize}
\item First item
%^--- 没有缩进
       \item    Second item with spaces
%     ^^^^^^^  缩进过多，且 item 后也有多余空格
  \item Third item
% ^^ 缩进 2 格
                \item Fourth item with too much indent
%               ^^^^^^^^^^^^^^^^ 缩进太多
  \end{itemize}

% ================================================================
%  问题 8：enumerate 列表同样缩进混乱
% ================================================================
\begin{enumerate}
    \item Step one
 %^^  缩进 4 格
 \item  Step two
%^ 缩进 1 格，且 Step 前多一个空格
\item Step three
%^ 没有缩进
                \item Step four
%               ^^^^^^^^^^^^^^^^ 缩进太多
  \end{enumerate}
% ^^ \end 缩进 2 格，与 \begin 不对齐

  \section{   Results and   Discussion   }
%         ^^^              ^^^  标题内部多余空格

We present our findings in Table~\ref{tab:results}.
% ~ 表示"不可断行空格"，确保"Table"和编号不会分到两行

% ================================================================
%  问题 9：表格中的列对齐不一致
%  \hline 有的缩进 2 格、有的顶格
%  数据行的 & 位置不统一，有的在前后有多余空格
% ================================================================
\begin{table}[h]
% [h] = 建议放在当前位置（here）
\centering
% 让表格居中
\begin{tabular}{lcr}
% {lcr} = 三列：左对齐 / 居中 / 右对齐
  \hline
  Method  & Accuracy & Time \\
  \hline
   A  &  95.2\%  &  1.3s \\
%  ^^ 列间距不统一
    B &   93.8\% &   2.1s \\
%  ^^^ 对齐方式与上一行不一致
   C  &   96.1\%  &   0.9s \\
  \hline
\end{tabular}
  \caption{  Experimental results  }
%         ^^                     ^^  caption 中多余空格
\label{tab:results}
\end{table}

  \subsection{Analysis}
% ^^ 多余缩进

  Our method outperforms baselines by a large margin.
  We attribute this improvement to our novel architecture.

\begin{equation}
  \text{Accuracy} = \frac{TP + TN}{TP + TN + FP + FN}
\end{equation}

  Another equation for reference:
  \begin{equation}
    \text{Loss} = -\sum_{i=1}^{n} y_i \log(\hat{y}_i)
  \end{equation}

  \section{Conclusion}
% ^^ 多余缩进

    In this paper, we have demonstrated the effectiveness of our approach.
    Future work includes extending the method to other domains.

% ================================================================
%  问题 10：thebibliography 缩进不一致
% ================================================================
  \begin{thebibliography}{9}
    \bibitem{ref1} Author A, ``Paper Title,'' Journal, 2024.
     \bibitem{ref2} Author B, ``Another Paper,'' Conference, 2023.
%    ^ 两个 \bibitem 缩进不一致
  \end{thebibliography}

\end{document}
